\chapter{Cases: Testing Reflexive Constraint Realism}
\label{ch:cases}

\begin{chapterguide}
This chapter tests \RCR\ across different sciences. Each case asks what was
measured, what constraints were established, which parts of the successful
representation were idealized, and what the case does not justify. The cases
include physics, biology, medicine, psychology, climate science, economics,
and social science.
\end{chapterguide}

A philosophy of science should be answerable to science as it is practiced.
The following cases are not offered as decisive proof of \RCR. They are stress
tests. Each case has a different pattern of access, evidence, intervention, and
institutional organization. If one framework can illuminate all of them
without pretending that their methods are identical, it has earned some
credibility.

\section{Physics: atoms, electrons, and theory change}

Nineteenth-century debates about atoms show why ontological commitment should
be graded. Kinetic and chemical theories used atoms to explain regularities,
but some scientists treated atoms as calculational devices rather than
established entities. Einstein's 1905 analysis of Brownian motion connected
observable particle movement with molecular-kinetic theory. Jean Perrin later
used several quantitative investigations of Brownian motion and related
phenomena to support the atomic theory \citep{einstein1905,perrin1913}.

The philosophical significance is not that one paper proved atoms in a
deductive sense. It is that independent relations converged: diffusion,
sedimentation, viscosity, and particle motion could be connected by a common
molecular framework. The framework generated quantitative constraints that
were not simply restatements of a single observation. Under \RCR, atoms move
from useful posit toward strong causal and structural commitment because the
same posit supports multiple measurement paths.

The electron provides a related case. J. J. Thomson's 1897 work on cathode
rays interpreted their deflection as evidence for negatively charged
constituents smaller than atoms. Subsequent measurements, including
Millikan's oil-drop work, connected charge and mass through distinct
experimental arrangements. The history is not a story of one experiment
silently revealing an entity. It is a sequence of instruments, calibrations,
controversies, and theoretical interpretations in which a stable causal role
emerged.

The Michelson--Morley experiment of 1887 illustrates a different point. It
did not by itself refute every version of an ether theory or uniquely establish
special relativity. Its result became significant within a larger network of
optics, mechanics, and later theoretical developments
\citep{michelsonmorley1887}. This is a case where the Duhem--Quine problem is
not an excuse for ignoring evidence. The experiment constrained a family of
models and contributed to a reorganization of theoretical structure.

\begin{casebox}
\textbf{\RCR\ verdict.} Physics supports realism about certain entities and
relations when they survive independent measurement and intervention, but it
does not support the assumption that every historical interpretation of those
entities was correct. What is retained may be a relation, a field equation, a
limiting case, or an interventionally stable capacity rather than a complete
ontology.
\end{casebox}

\section{Biology: the structure of DNA}

The 1953 proposal by James Watson and Francis Crick represented DNA as a
double helix with paired bases \citep{watsoncrick1953}. The result was not
deduced from a single image. It combined chemical constraints, model-building,
and X-ray diffraction evidence, including work by Rosalind Franklin and
Raymond Gosling. The published structure was a selective representation: it
captured a geometry and pairing relation without being a complete theory of
replication, gene regulation, or development.

This case displays all four layers of scientific success.

\begin{enumerate}
\item \textbf{Correspondence:} DNA has a molecular structure that constrains
its chemical behaviour.
\item \textbf{Structural accuracy:} the double-helical arrangement preserves
relations among the backbone and base pairs.
\item \textbf{Interventional adequacy:} later molecular biology could alter,
sequence, and manipulate nucleic-acid systems in ways that depended on the
structure.
\item \textbf{Pragmatic adequacy:} the representation supported new research
questions and technologies.
\end{enumerate}

The case also shows why social epistemology matters. The historical credit
attached to a discovery does not perfectly track the distribution of data,
labour, and interpretation. A theory can be scientifically strong while the
institutional history of recognition is unjust. Epistemic evaluation should
therefore separate the truth of a structural claim from the fairness of the
community that produced and rewarded it.

\section{Medicine: the MRC streptomycin trial}

The Medical Research Council's 1948 controlled investigation of streptomycin
for pulmonary tuberculosis is an important case of experimental and
institutional design \citep{mrc1948}. The trial used allocation procedures
intended to prevent foreknowledge of treatment assignment and compared
outcomes under a defined protocol. The result did not establish that
streptomycin was a complete solution to tuberculosis. Resistance, toxicity,
disease severity, and combination therapy remained relevant.

Its philosophical importance lies in the separation of causal question from
clinical enthusiasm. The question was not merely whether patients receiving
streptomycin appeared to improve. It was whether outcomes differed under a
controlled comparison that reduced selection bias. Randomization did not make
the study theory-free. It required definitions of disease, treatment, outcome,
follow-up, and missingness. But it created a constraint that a persuasive
narrative alone could not easily evade.

Under \RCR, the trial strongly supports a scoped causal claim: under the
studied conditions, the treatment affected specified outcomes in the studied
population. It does not license an unqualified claim that the drug works in
all settings, for all forms of disease, or as a universal mechanism. The
intervention boundary matters.

\section{Psychology: replication and context}

The Open Science Collaboration's 2015 project attempted to replicate 100
psychological studies and reported that replicated effects were, on average,
smaller than the original effects \citep{osc2015}. The project is often
summarized as a simple failure rate. That summary is inadequate. Replication
success can be evaluated by significance, effect size, interval overlap,
prediction intervals, or theoretical fit. Each criterion answers a different
question.

The case illustrates three \RCR\ principles.

First, a result is not fully objective because it passed one statistical
threshold. Publication selection, flexible analysis, and measurement can
inflate apparent evidence.

Second, a failed replication is evidence about the claim and the conditions,
not a metaphysical verdict on a whole discipline. If an effect depends on
population, context, stimulus, or implementation, a conceptual replication
may reveal a boundary rather than a simple falsehood.

Third, social organization is part of epistemology. A field that rewards
novelty and treats correction as failure weakens its own correction capacity.
Preregistration, open materials, independent analyses, and better incentives
can improve objectivity without guaranteeing uniform results.

\section{Climate science: convergence across models and evidence}

Climate science is a large-scale case of model-mediated, multi-source
inference. No single thermometer, simulation, or statistical method establishes
the whole climate system. The evidence includes instrumental records,
paleoclimate proxies, physical theory, radiative measurements, ocean heat
content, cryosphere changes, attribution studies, and model projections.

The IPCC's Sixth Assessment Synthesis Report, released in 2023, assesses that
human influence has warmed the atmosphere, ocean, and land, and it expresses
the assessment with calibrated language about confidence and likelihood
\citep{ipcc2023}. The report does not remove uncertainty. It distinguishes
established observed changes, attribution, projected risks, and the conditions
under which projections vary.

Climate science therefore tests \RCR\ in a domain where direct intervention on
the global system is unavailable and the target is dynamic. Its objectivity
comes from convergence among independent constraints and from the physical
mechanisms linking greenhouse-gas concentrations, energy balance, circulation,
and observed changes. Models are not literal replicas of the atmosphere. They
are idealized tools whose reliability depends on hindcasts, process validation,
inter-model comparison, and observation.

The case also shows why values enter without deciding the physical conclusion.
Choices about acceptable risk and vulnerable populations affect policy
interpretation, but they do not determine whether human forcing has altered
the climate system.

\section{Economics: policy models and the Lucas critique}

Economic models demonstrate the importance of reflexivity. A policy can change
the expectations and behaviour that generated an earlier statistical
relationship. Robert Lucas's critique argued that policy evaluation based on
historical correlations can fail when policy changes the decision rules of
agents \citep{lucas1976}.

The philosophical lesson is not that economic models are unscientific. It is
that transportability requires a mechanism. A model estimated under one policy
regime may not predict under another if households, firms, or institutions
adapt. Causal identification must therefore address expectations,
equilibrium, institutional rules, and the time path of adjustment.

\RCR\ recommends separating several claims:

\begin{itemize}
\item a historical association among variables;
\item a structural model of decision rules;
\item a causal effect under a specified intervention;
\item a forecast conditional on agents and institutions not changing in
unmodelled ways.
\end{itemize}

An economic model can be useful and structurally informative without literally
describing every agent. Its objectivity depends on the stability of the
mechanism under the policy intervention being considered.

\section{Social science: field experiments and discrimination}

Field experiments can create causal constraints in settings where laboratory
control is limited. Bertrand and Mullainathan's audit study sent matched
resumes with different names to employers and examined differences in
callbacks \citep{bertrandmullainathan2004}. The design was intended to isolate
the effect of perceived racial identity on employer response while keeping
other resume information comparable.

The study does not identify every mechanism of discrimination. Names may signal
more than one social perception; employers may differ in interpretation; and
callback differences do not exhaust later hiring processes. It does, however,
show how a social category can be investigated through an intervention that
changes a signal while holding much of the application constant.

The case supports a scoped claim about employer responses under the study's
conditions. It also demonstrates the reflexive nature of social research:
categories such as race are socially produced, but their effects on decisions
can be objectively studied. The reality of a social mechanism is not reduced
to a biological essence or dissolved into language.

\section{Cross-case comparison}

\begin{table}[ht]
\centering
\smalltable
\begin{tabularx}{\textwidth}{p{0.20\textwidth}p{0.22\textwidth}Y Y}
\toprule
Field & Main access & Strongest constraint & Main limit\\
\midrule
Physics & Instrumental interaction and theory & Convergent measurement and intervention & Theory change can alter ontology\\
Biology & Structure, experiment, mechanism & Cross-scale causal integration & Complex context and multiple levels\\
Medicine & Controlled intervention & Randomized treatment comparison & Scope, ethics, heterogeneity\\
Psychology & Experiment, measurement, replication & Independent reproduction and design transparency & Context and measurement variation\\
Climate & Observation, physical model, attribution & Convergence across sources and scales & No global intervention; projection uncertainty\\
Economics & Historical data, models, policy variation & Mechanism under changing expectations & Reflexivity and transportability\\
Social science & Field experiment, records, interpretation & Scoped intervention and plural critique & Categories and institutions can change\\
\bottomrule
\end{tabularx}
\caption{Different sciences realize objectivity through different constraint
profiles.}
\label{tab:case-comparison}
\end{table}

The cases support a modest generalization. Objective science is not defined by
one preferred method. It is defined by the capacity of a method and institution
to create relevant constraints, expose error, and state what the result does
not establish.

\section*{Chapter summary}

Physics shows how entities and structures earn commitment across independent
measurements. Biology shows how model-building and social credit can diverge.
Medicine shows the causal value of controlled comparison and the importance of
scope. Psychology shows why replication is a graded and contextual test.
Climate science shows how convergence can support strong conclusions without
eliminating uncertainty. Economics shows that policy changes can alter the
mechanisms being estimated. Social science shows that socially constructed
categories can have objectively testable causal effects.

\begin{discussion}
\begin{enumerate}
\item Which case provides the strongest evidence for entity realism?
\item Which case most clearly separates pragmatic usefulness from causal
adequacy?
\item How should a climate or economic model report uncertainty that comes from
future institutional change?
\end{enumerate}
\end{discussion}

\begin{furtherreading}
Read the primary studies cited in the cases. For broader comparisons, consult
Hacking (1983), Woodward (2003), Cartwright (1999), Craver (2007), and Kitcher
(2001).
\end{furtherreading}

